\section{结论}
\label{sec:conclusion}

本文区分上游 ARA 原型、本地 point-v1 和未有实测结果的 trajectory-v2，保留历史硬近邻目标推导，并形式化本地缓存输出增量、软邻域、非负间隔惩罚、完整因子重置与规范化检查。任意秩仅限定上游全矩阵参数化不显式约束更新秩；本地实验采用秩至多 128 的因子。软目标的零下界与数值保护都不能建立全局收敛或机制因果性。

归档 point-v1 的 120 个试次全部完成，但没有候选通过联合效果门槛。最低 Keywords 为 0.54，对应首词元 KL 为 0.089975，验收状态为 failed，没有通过验收的适配器。工程完成与效果失败同时成立。单次自适应验证搜索、逐提示证据缺失及运行版本记录不完整进一步限制了结论强度，不能由这些代理确认语义拒答下降、通用能力保持或跨方法优势。

trajectory-v2 已实现多位置轨迹、连续评分和更完整的验收控制，其收益仍待实验检验。后续应在固定模型、隔离数据、统一预算和独立语义审计下，检验激活漂移、较晚词元覆盖与搜索边界假设，再讨论结构变化是否带来可重复的行为改进。
